\section{Introduction}
\label{sec:introduction}

Large language models (LLMs) have demonstrated strong capabilities in general-purpose code generation \cite{vaswani2017attention, openai2023gpt4, yang2025qwen3}. However, applying LLMs to EDA scripting—where engineers routinely manipulate in-memory design databases, create layout instances, route schematics, and query netlists—remains challenging. First, EDA APIs are long-tailed and tool-specific; many rarely appear in public training corpora. Second, correctness cannot be verified statically: it depends on observable side effects in a real EDA environment, such as modified layouts or updated schematics. Third, without execution feedback, static generation cannot diagnose errors or iteratively refine incorrect code \cite{chen2024dawning, blocklove2024eda}. Existing LLM-based EDA work has mainly focused on HDL generation \cite{thakur2023verigen, liu2023rtlcoder}, while scripting with real tool APIs—especially with incomplete documentation and sandbox execution—remains underexplored.
% 大语言模型在通用代码生成方面表现强大，但应用于EDA脚本——工程师日常操作内存中的设计数据库、创建版图实例、布线原理图、查询网表——仍面临挑战。第一，EDA API具有长尾性和工具特定性，许多很少出现在公开训练语料中。第二，正确性无法静态验证：它取决于真实EDA环境中的可观察副作用，如版图修改或原理图更新。第三，没有执行反馈，静态生成无法诊断错误或迭代修正代码。现有LLM-for-EDA工作主要集中在HDL生成，而面向真实工具API、沙盒执行和不完整文档的脚本生成仍研究不足。

This paper presents \textbf{ZhuLong}, an execution-grounded LLM coding agent for PyAether and SKILL. Named after a mythical creature whose opened eyes bring light, ZhuLong aims to illuminate opaque EDA APIs and automate tool-specific scripting workflows. Built on Cline-CLI \cite{cline2024cline}, ZhuLong extends three MCP tools \cite{lee2025autoeda}—\texttt{search\_apis}, \texttt{get\_api\_details}, and \texttt{run\_code}—that form a closed loop: retrieve candidate APIs, inspect their documentation, execute code in the EDA environment with execution feedback, and iteratively refine based on observed errors or outputs. All tools accept a unified \texttt{lang} parameter, supporting both PyAether and SKILL with identical agent logic. To further address incomplete or outdated documentation, ZhuLong introduces an offline \textbf{API self-exploration mechanism} that proactively explores undocumented API behaviors in the sandbox via \textbf{counterfactual experimentation}—deliberately testing alternative parameter values, observing execution outcomes, inferring constraints and usage patterns, and storing the discovered information as enhanced API documentation for future queries.
% 本文提出\textbf{ZhuLong}，一个面向PyAether和SKILL的执行驱动LLM代码智能体。其名源自中国神话中睁眼带来光明的神兽，寓意照亮不透明的EDA API并自动化工具特定的脚本流程。ZhuLong基于Cline-CLI构建，扩展三个MCP工具——\texttt{search\_apis}、\texttt{get\_api\_details}和\texttt{run\_code}——形成闭环：检索候选API、查阅文档、在沙盒EDA环境中执行代码、基于观察到的错误或输出迭代修正。所有工具接受统一的\texttt{lang}参数，以相同智能体逻辑同时支持PyAether和SKILL。为了进一步解决文档不完整或过时的问题，ZhuLong引入了离线\textbf{API自探索机制}，通过沙盒中的\textbf{反事实实验}主动探索未文档化的API行为——有意识地测试不同参数取值、观察执行结果、推断约束和使用模式、将发现的信息存储为增强的API文档供后续查询使用。

The main contributions are:
\begin{itemize}
    \item \textbf{First execution-grounded agent for commercial EDA scripting.} We present the first system that validates execution feedback in PyAether and SKILL via sandbox execution, delivering a +43 pp accuracy gain over static RAG.
    % 首个面向商业EDA脚本的执行驱动智能体。我们提出了首个通过沙盒执行验证PyAether和SKILL执行反馈的系统，相比静态RAG带来+43个百分点的准确率提升。
    \item \textbf{Offline API self-exploration for EDA.} A complementary mechanism that contributes +3.2 pp accuracy gain and reduces the average number of tool calls per trace by 22.1\% through counterfactual experimentation on undocumented APIs. To our knowledge, this is the first mechanism that proactively discovers and codifies undocumented API behaviors in the EDA domain.
    % \textbf{面向EDA的离线API自探索。}一种互补机制，通过对未文档化API进行反事实实验，贡献了+3.2个百分点的准确率提升，并将每条trace的平均工具调用次数减少22.1\%。据我们所知，这是首个在EDA领域中主动发现并系统记录未文档化API行为的机制。
    \item \textbf{EDA-Eval-PyAether: the first public benchmark for PyAether scripting.} A benchmark of 158 real-world tasks with assertion-based execution, enabling rigorous evaluation of LLM-generated EDA scripts. The complete ZhuLong system achieves 78.5\% Pass@1 on this benchmark.
    % EDA-Eval-PyAether：首个PyAether脚本公开评测集。包含158个真实任务的基准，采用基于断言的执行验证，支持对LLM生成的EDA脚本进行严格评估。完整ZhuLong系统在该基准上达到78.5\%的Pass@1。
\end{itemize}